\section{机器人、信息市场与物理世界}

\subsection{原子比比特难一百万倍}

Karpathy 从 Tesla 的自动驾驶经历出发：十年前有大量自动驾驶创业公司，大多数最终没活下来。原因是资本密集、时间漫长、需要巨大的信念感。

他重申了数字优先的路线图：数字空间的 unhobbling $\rightarrow$ 数字/物理接口 $\rightarrow$ 物理世界全面渗透。物理世界的总可寻址市场可能远大于数字空间，但难度也大一百万倍。

\begin{knowledgebox}{Periodic Labs}
由前 OpenAI 研究 VP Liam Fedus 和前 DeepMind 材料科学主管 Ekin Dogus Cubuk 联合创办，2025年获得3亿美元种子轮融资。公司目标是构建能自主运行物理实验的"AI 科学家"。这是数字/物理接口的典型案例——传感器不只是摄像头，还有昂贵的实验室设备。
\end{knowledgebox}

\subsection{缺失的信息市场}

Karpathy 指出了一个有意思的缺失：如果 Polymarket 等预测市场有越来越多的自主 Agent 参与，如果伊朗正在发生什么事，从德黑兰拍一张照片应该值10美元——不是人在看，是 Agent 在试图判断预测市场和股票的走势。但目前没有这种机制。

他引用了 Daniel Suarez 的科幻小说《Daemon》——书中的超级智能把人类既当传感器又当执行器，社会围绕着这台机器的需求重新组织。他觉得某种类似的事情正在发生：越来越多的自动化有特定需求，人类会开始服务于那台机器的需求。

\subsection{本章小结}

物理世界是更大的市场但更难攻克。数字/物理接口（传感器和执行器）是近期最有趣的创业方向。信息市场的缺失意味着 Agent 经济的基础设施还远未建成。


